\documentclass[10pt]{amsart}
\usepackage[a4paper,margin=23mm]{geometry}
\usepackage{amsmath,amssymb,amsthm,mathtools,hyperref,graphicx,xspace,letltxmacro,xpatch}
\usepackage[english]{babel}
\newtheorem{theo}{Theorem}[section]
\newtheorem{lemm}[theo]{Lemma}\newtheorem{prop}[theo]{Proposition}\newtheorem{coro}[theo]{Corollary}\newtheorem{conj}[theo]{Conjecture}
\theoremstyle{definition}\newtheorem{defi}[theo]{Definition}\newtheorem{rema}[theo]{Remark}\newtheorem{exam}[theo]{Example}\newtheorem{exem}[theo]{Example}\newtheorem{nota}[theo]{Notation}
\newcommand{\ie}{i.e.}\newcommand{\eg}{e.g.}\newcommand{\resp}{resp.}\newcommand{\skpt}{\par\smallskip\noindent}
\emergencystretch=3em\pagestyle{plain}
%~Mouliné par MaN_auto v.0.23.0 2020-05-14 15:26:36


\DeclareMathOperator{\vol}{Vol}
\DeclareMathOperator{\conv}{conv}

\newcommand*{\dd}{\mathrm{d}}


\newcommand{\wet}[1]{w(#1)}
\newcommand{\eps}{\varepsilon}
\newcommand{\Nbb}{\mathbb{N}}%%\Nbb
\newcommand{\Rbb}{\mathbb{R}}%\Rbb
\newcommand{\Ebb}{\mathbb{E}}%\Ebb
\newcommand{\Hcal}{\mathcal{H}} %%\Hcal


%%désactivés pour l'instant
\newcommand{\pth}[1]{\left(#1 \right)}
\newcommand{\ptb}[1]{\left[#1 \right]}
\newcommand{\blue}[1]{{\color{blue}#1}}


%%\AtBeginDocument{
%\newtheorem{myname}[cdrthm]{My beautiful theorem} }


\newcommand*{\mtfrac}[2]{\frac{#1}{#2}}% si jamais l'auteur veut absolument des tfrac


%%%%%%%%%%%%%%%%%%%%%%%%%%%%%%%%%%%%%%%%%%%%%%%%%%%%

\graphicspath{{./figures/}}

\newcommand*{\mk}{\mkern -1mu}
\newcommand*{\Mk}{\mkern -2mu}
\newcommand*{\mK}{\mkern 1mu}
\newcommand*{\MK}{\mkern 2mu}

\hypersetup{urlcolor=purple, linkcolor=blue, citecolor=red}

\newcommand*{\romanenumi}{\renewcommand*{\theenumi}{\roman{enumi}}}
\newcommand*{\Romanenumi}{\renewcommand*{\theenumi}{\Roman{enumi}}}
\newcommand*{\alphenumi}{\renewcommand*{\theenumi}{\alph{enumi}}}
\newcommand*{\Alphenumi}{\renewcommand*{\theenumi}{\Alph{enumi}}}
%%%%%%%%%%%%%%%%%%%%%%%%%%%%%%%%%%%%%%%%%%%%%%%%%%%%%%%%%

